% =====================================================================
%  PLANTILLA — Slides académicas de física (estilo de S. Gaviria Giraldo)
%  Derivada de baseline/slides.tex. Copiar este archivo como punto de
%  partida y reemplazar el contenido de los frames de ejemplo.
%  Compilar:  tectonic slides.tex   (o bien  pdflatex slides.tex)
% =====================================================================
\documentclass[aspectratio=169,11pt]{beamer}

\usepackage{iftex}
\ifPDFTeX
  \usepackage[T1]{fontenc}
  \usepackage[utf8]{inputenc}
  \usepackage{lmodern}
\fi
\usepackage[spanish,es-nodecimaldot,es-noshorthands]{babel}
\usepackage{amsmath,amssymb}
\usepackage{graphicx}
\usepackage{booktabs}
\usepackage{array}
\usepackage{tikz}
\usetikzlibrary{arrows.meta,positioning,calc}
\graphicspath{{figures/}}

% ---------------------------------------------------------------------
%  Estilo sobrio
% ---------------------------------------------------------------------
\definecolor{azul}{RGB}{20,83,56}         % color principal (verde)
\definecolor{azulclaro}{RGB}{228,241,233} % verde claro
\definecolor{gris}{RGB}{90,90,90}
\definecolor{acento}{RGB}{176,58,46}

\usefonttheme{professionalfonts}
\setbeamertemplate{navigation symbols}{}
\setbeamercolor{structure}{fg=azul}
\setbeamercolor{frametitle}{fg=white,bg=azul}
\setbeamercolor{framesubtitle}{fg=azul,bg=azulclaro}
\setbeamercolor{framenumber}{fg=azulclaro,bg=azul}
\setbeamercolor{title}{fg=white,bg=azul}
\setbeamercolor{block title}{fg=white,bg=azul}
\setbeamercolor{block body}{fg=black,bg=azulclaro}
\setbeamertemplate{blocks}[rounded][shadow=false]
\setbeamertemplate{itemize items}{\small$\blacktriangleright$}
\setbeamerfont{frametitle}{size=\large,series=\bfseries}
\setbeamerfont{framesubtitle}{size=\small,series=\mdseries,shape=\itshape}
\setbeamerfont{framenumber}{size=\footnotesize,series=\bfseries}
\setbeamerfont{caption}{size=\scriptsize}
\setbeamertemplate{caption}{\raggedright\insertcaption\par}

% Título de cada slide en barra redondeada (con número de slide);
% subtítulo en caja clara con filete lateral.
\makeatletter
\setbeamertemplate{frametitle}{%
  \vspace*{0.5em}%
  \begin{beamercolorbox}[wd=\textwidth,sep=0.45em,leftskip=0.3em,rightskip=0.3em,rounded=true,shadow=false]{frametitle}%
    \usebeamerfont{frametitle}\strut\insertframetitle\hfill
    {\usebeamerfont{framenumber}\usebeamercolor[fg]{framenumber}\insertframenumber\,/\,\inserttotalframenumber}%
  \end{beamercolorbox}%
  \ifx\insertframesubtitle\@empty\else
    \par\vspace{0.3em}%
    \hbox{%
      {\color{acento}\vrule width 2.5pt}%
      \begin{beamercolorbox}[wd=\dimexpr\textwidth-2.5pt\relax,sep=0.3em,leftskip=0.3em]{framesubtitle}%
        \usebeamerfont{framesubtitle}\mdseries\strut\insertframesubtitle
      \end{beamercolorbox}}%
  \fi
  \vspace{0.2em}%
}
\makeatother

% Sin pie de página
\setbeamertemplate{footline}{}

\newcommand{\kB}{k_{B}}
\newcommand{\fig}[1]{{\scriptsize\color{gris}#1}}

% Estilos TikZ reutilizables (cabecera verde + caja clara + flecha roja)
\tikzset{
  cab/.style={fill=azul,text=white,rounded corners=3pt,font=\small\bfseries,
              text width=3.8cm,align=center,inner sep=3pt,minimum height=0.6cm},
  caja/.style={draw=azul,fill=azulclaro,rounded corners=3pt,font=\footnotesize,
               text width=3.8cm,align=flush left,inner sep=5pt,anchor=north},
  paso/.style={-{Stealth},very thick,acento},
  etapa/.style={draw=azul,fill=azulclaro,rounded corners=3pt,font=\footnotesize,
                align=center,inner sep=4pt},   % caja de un paso en cadenas/algoritmos
}

% ---------------------------------------------------------------------
\title{Título del trabajo}
\author{Autor}
% =====================================================================

\begin{document}

% ---------------------------------------------------------------- portada
\begin{frame}[plain]
  \vfill
  \begin{beamercolorbox}[wd=\textwidth,sep=1em,center,rounded=true,shadow=false]{title}
    {\Large\bfseries Título principal del trabajo:\\[0.3em]
     segunda línea del título\par}
    \vspace{0.5em}
    {\color{white!60!azul}\rule{0.45\linewidth}{0.4pt}\par}
    \vspace{0.35em}
    {\normalsize\itshape Nombre del evento -- Ciudad, año\par}
  \end{beamercolorbox}
  \vspace{1.2em}
  \centering
  {\large Autor~Uno,\; Autor~Dos\par}
  \vspace{0.5em}
  {\small\color{gris}Instituto / Departamento\\ Universidad, Ciudad, País\par}
  \vspace{0.7em}
  \makebox[\textwidth][c]{\begin{minipage}{0.62\textwidth}
  \begin{beamercolorbox}[wd=\textwidth,sep=0.45em,center,rounded=true]{framesubtitle}
    \small Dato de contexto (revista, tesis, proyecto)\\
    \footnotesize línea secundaria (doi, repositorio, financiación)
  \end{beamercolorbox}
  \end{minipage}}\par
  \vfill
\end{frame}

% ---------------------------------------------------------------- definición + figura/tabla + dos bloques
\begin{frame}{¿Qué es el concepto X?}
  \framesubtitle{\textbf{Definición:} una frase precisa que el público pueda leer en 5 segundos.}
  \begin{columns}[c,onlytextwidth]
    \begin{column}{0.42\textwidth}
      \centering
      \begin{tikzpicture}[x=0.82cm,y=0.62cm,>=stealth]
        \draw[thick,azul] plot[smooth,tension=0.7] coordinates
          {(0,2.1) (0.8,0.45) (1.6,1.75) (2.6,0.8) (3.5,1.95) (4.4,0.3) (5.3,1.9)};
        \fill[acento] (0.8,0.62) circle (0.11);
        \draw[->,thick,acento] (0.95,0.75) .. controls (1.2,2.6) and (2.2,2.4) .. (2.45,1.0);
        \node[font=\scriptsize] at (2.9,-0.55) {coordenada};
      \end{tikzpicture}\\[0.1em]
      {\small $r = \nu_0 \exp\!\left(-\dfrac{E}{\kB T}\right)$}
    \end{column}
    \begin{column}{0.55\textwidth}
      \centering\footnotesize
      \renewcommand{\arraystretch}{1.15}
      \begin{tabular}{@{}lll@{}}
        \toprule
        Método & Escala & Qué entrega \\
        \midrule
        A & fs & algo \\
        \textbf{B} & \textbf{s} & \textbf{lo nuestro} \\
        \bottomrule
      \end{tabular}
    \end{column}
  \end{columns}
  \vspace{0.5em}
  \begin{columns}[T,onlytextwidth]
    \begin{column}{0.485\textwidth}
      \begin{block}{Idea 1}
        \footnotesize Frase corta $\Rightarrow$ \emph{consecuencia}.
      \end{block}
    \end{column}
    \begin{column}{0.485\textwidth}
      \begin{block}{Idea 2}
        \footnotesize Frase corta $\Rightarrow$ \emph{consecuencia}.
      \end{block}
    \end{column}
  \end{columns}
\end{frame}

% ---------------------------------------------------------------- ecuaciones + figura
\begin{frame}{Ecuaciones con figura}
  \begin{columns}[c,onlytextwidth]
    \begin{column}{0.5\textwidth}
      \begin{align}
        \Delta\mu &= \kB T \ln(\sigma + 1) \tag{1}\\[0.3em]
        \Delta G  &= -N\Delta\mu + \gamma A \tag{2}
      \end{align}
      \begin{block}{}
        \begin{equation}
          r = K_0\exp\!\left(-\frac{\Delta G}{\kB T}\right) \tag{3}
        \end{equation}
      \end{block}
      {\footnotesize $\gamma$: energía de borde; \ $A$: área nueva}
    \end{column}
    \begin{column}{0.46\textwidth}
      \centering
      \fbox{\parbox[c][3.2cm][c]{0.9\linewidth}{\centering\color{gris}figura}}\\
      \fig{Fig.~1. Descripción breve de la figura.}
    \end{column}
  \end{columns}
\end{frame}

% ---------------------------------------------------------------- flujo de cajas TikZ
\begin{frame}{Motivación y problema}
  \framesubtitle{Mensaje central de la slide en una línea.}
  \centering
  \begin{tikzpicture}
    \node[cab] (h1) at (0,0) {Motivación};
    \node[caja,minimum height=2.4cm] (b1) at ([yshift=-2pt]h1.south) {\textbf{Punto} clave\\[0.4em] otro punto};
    \node[cab] (h2) at (4.9,0) {Causa};
    \node[caja,minimum height=2.4cm] (b2) at ([yshift=-2pt]h2.south) {\textbf{Punto} clave};
    \node[cab] (h3) at (9.8,0) {Problema};
    \node[caja,minimum height=2.4cm] (b3) at ([yshift=-2pt]h3.south) {\textbf{Punto} clave};
    \draw[paso] (b1.east) -- (b2.west);
    \draw[paso] (b2.east) -- (b3.west);
    \node[fill=azul,text=white,rounded corners=3pt,font=\small,text width=13.3cm,
          align=center,inner sep=5pt,anchor=north] (prop) at ([yshift=-0.45cm]b2.south)
          {\textbf{Propuesta:} una línea con la idea del trabajo};
    \draw[paso] (b3.south) -- (b3.south |- prop.north);
  \end{tikzpicture}
\end{frame}

% ---------------------------------------------------------------- conclusiones
\begin{frame}{Conclusiones}
  \begin{columns}[T,onlytextwidth]
    \begin{column}{0.56\textwidth}
      \begin{itemize}\setlength{\itemsep}{0.7em}
        \item Resultado con \textbf{palabra clave}
        \item Resultado cuantitativo: $R^2 = 0.96$
      \end{itemize}
    \end{column}
    \begin{column}{0.40\textwidth}
      \begin{block}{Limitaciones}
        \small Frase corta
      \end{block}
      \vspace{0.4em}
      \begin{block}{Perspectivas}
        \small Frase corta
      \end{block}
    \end{column}
  \end{columns}
\end{frame}

% ---------------------------------------------------------------- cierre
\begin{frame}{¡Gracias por su atención!}
  \vfill
  \begin{block}{}
    \centering
    \Large ¿Preguntas?
  \end{block}
  \vfill
\end{frame}

\end{document}
